% =====================================================================================
\chapter{METHODOLOGY: WELDSET, A WELD SEAM DATASET WITH ANALYTIC GROUND TRUTH}
\label{chp:method_weldset}
% Second shape 2026-10-20, rewritten after a line-by-line review of weld_generator
% (thesis/notes/review_2026-10-20/C*.md). Every value below is the code's as of 2026-10-20
% (branch tack_welding); ranges are "as coded" and, where the corpus differs, "as emitted".
% Numbers produced by the benchmark (corpus sizes, scores, gate statistics) live on the
% desktop and belong to Chapter 6: they are tagged DESKTOP.
% Known defects found in the review: thesis/notes/weldset_code_review_2026-10-20.md.
% =====================================================================================

Methods that extract weld seams from point clouds are evaluated against labels that people
drew by hand, on data that is rarely released, and the labels themselves are less accurate
than the accuracies claimed. In this chapter, we describe WeldSet, a generator of weld
assemblies in which the seam is constructed rather than labelled, so that it is known
exactly. Section~\ref{sec:ws_goals} states the motivation and the design goals and
Section~\ref{sec:ws_pipeline} gives the generation pipeline as a whole. The sections after
them follow the pipeline: parts and joints (\S\ref{sec:ws_parts}), curved seams
(\S\ref{sec:ws_curves}), the seam rule (\S\ref{sec:ws_d4}), the tack rule
(\S\ref{sec:ws_tacks}), the most probable seam (\S\ref{sec:ws_mps}) and the sensor model
(\S\ref{sec:ws_sensor}). The corpus (\S\ref{sec:ws_corpus}) and the evaluation protocol
with its seven re-implemented methods (\S\ref{sec:ws_protocol}) come next, then the two seam
extractors proposed in this work (\S\ref{sec:ws_own}), the link between WeldSet and the
robotic cell of Chapter~\ref{chp:method_cell} (\S\ref{sec:ws_link}), and the
annotation-error experiment (\S\ref{sec:ws_annotation}).

% -------------------------------------------------------------------------------------
\section{Motivation and Design Goals}
\label{sec:ws_goals}
Three faults are common to the evaluation of seam-extraction methods. First, the reference
is inexact: hand-annotated or estimated seams carry errors of the same order as the errors
reported for the methods, so a sub-millimetre accuracy cannot be resolved against them.
Second, the conditions are not varied: thickness, gap, angle, sensor noise and viewpoint are
fixed by the parts that happened to be scanned, so a method's dependence on them cannot be
separated. Third, the data are seldom available, which rules out an independent comparison.

In this thesis, we follow one rule that removes the first fault and makes the other two
cheap to address: \emph{ground truth is constructed, never detected}. The generator samples
a joint, places the parts by closed-form transforms and computes the seam from the same
transforms, exact to floating point. No annotator and no estimator is in the label path.
Four design habits follow from it:
\begin{itemize}
\item \textbf{Truth is separated from convention.} Everything that is a convention rather
  than geometry (the tack layout, the most probable seam, the noise realisation) is a
  versioned rule applied to the stored scene, so it can be re-applied when the rule changes
  without regenerating the geometry.
\item \textbf{Exactness has a contract.} Curves are analytic with analytic positions; parts
  are exact primitives; meshes used for ray casting have a recorded chord error under a
  fixed budget.
\item \textbf{Determinism by content hash.} A scene is identified by its configuration and
  seed and verified by a hash of its contents, under a pinned software environment.
\item \textbf{Accessibility.} The first tier needs no renderer, simulator or GPU; it uses
  NumPy and trimesh only.
\end{itemize}
\TODO{Three-motivations paragraph in Kurt's style (Pattern 14, notes/writing\_style\_guide.md)
for why synthetic, analytic data: labour, ground truth, controllability.}
\SRC{weld\_generator MANUAL.md \S1--2; dataset\_plan.md \S1 (D1--D27); README.MD.}

% -------------------------------------------------------------------------------------
\section{The Generation Pipeline}
\label{sec:ws_pipeline}
A scene is generated by Algorithm~\ref{alg:ws_generate}. The pipeline is inverted with
respect to perception: instead of finding a seam in a scan, it starts from the joint and
derives everything else, so that the seam is never solved for.

\begin{algorithm}[htbp]
\caption{Generating one WeldSet scene}
\label{alg:ws_generate}
\begin{algorithmic}[1]
\REQUIRE configuration $c$, seed $n$
\ENSURE scene description, point cloud, seam and tack arrays, or a recorded rejection
\STATE spawn eight named random substreams from $n$ (geometry, defects, seam curve, placement, sampling, camera, \dots)
\STATE sample the joint: type, part dimensions, thicknesses, included angle; and the defects: gap, misalignments, quality target
\STATE build the parts as exact primitives and place them by closed-form transforms (curved families: draw the seam curve first, then derive the parts from it)
\STATE enumerate the candidate seams and judge each by the accessibility rule (Algorithm~\ref{alg:ws_access}); keep the weldable ones as ground truth
\STATE draw a camera pose around the joint (\S\ref{sec:ws_sensor})
\STATE sample every face area-uniformly at the drawn density; store exact normals and part and face identifiers
\STATE compute per point: visibility from the camera, and exterior (reachable by a scan)
\IF{no weldable seam is visible enough}
  \STATE record the seed as skipped with its reason and \textbf{return}
\ENDIF
\STATE apply the rule blocks: tacks (Algorithm~\ref{alg:ws_tacks}) and the most probable seam
\STATE write the scene, its arrays and its content hash
\end{algorithmic}
\end{algorithm}

Seeds whose scene shows no weldable seam are not replaced: they are recorded with the
reason, so that the selection of a corpus is a pure function of the configuration and the
seed range. Because rejections depend on the joint type, the yield (emitted scenes per
attempted seed) differs between types, and it is reported with every per-type result.
\TODO{Yields per type (README, Phase 3: T 79\,\%, butt 57\,\%, corner 44\,\%, edge 42\,\%, lap
36\,\%) and per stratum of the final corpus (DESKTOP).}
\SRC{weldgen/scene.py, scene\_curved.py, rng.py, hashing.py, writer.py; README.MD.}

\paragraph{Determinism.} Random draws are taken from eight named substreams spawned from the
seed, and new draws are appended, never inserted, so that adding a parameter does not change
existing scenes. A scene identifier is the configuration identifier (a hash of the resolved
configuration) and the seed; its content hash covers the canonical description and every
array. Hashes are reproducible only under the pinned environment: on 2026-09-11 every hash
changed between two NumPy releases although the arrays agreed to $10^{-5}$ relative.
\TODO{State the pinned versions (NumPy, trimesh, Python) in Chapter 5 and pin
requirements.txt. Note: the noise and render substreams are named but not used (the noise is
seeded by the scene seed): word this correctly or fix the code (weldset\_code\_review, P3).}

% -------------------------------------------------------------------------------------
\section{Parts, Joints and Assemblies}
\label{sec:ws_parts}
Every part is an exact primitive with a face registry, so that every face of every part has
a name, an analytic plane or surface and an outward normal (Table~\ref{tab:ws_primitives}).
The thickness axis of a plate is always its local $w$ axis, and its two broad faces are
$\pm w$.
\begin{table}[htbp]
\caption{WeldSet primitives.}
\label{tab:ws_primitives}
\centering\small
\begin{tabular}{@{}L{0.2\textwidth}L{0.42\textwidth}L{0.3\textwidth}@{}}
\toprule
Primitive & Parameters & Models \\
\midrule
Slab & dimensions $(L,W,t)$, pose & a rectangular plate; the fixture \\
Prism & convex outline, thickness, pose & a plate with a polygonal outline \\
PreparedSlab & length, width, thickness, ISO 9692-1 groove (kind, bevel, root face, radius) & a plate with a grooved edge \\
PreparedPrism & outline, thickness, groove & both of the above \\
Tube & outer radius, wall, length, base cut (none, a plane, or a cylinder), gap & a pipe with a straight, mitred or saddle-cut end \\
SweptSlab & a planar spine curve, offsets, height & a band along a curve: a box section, a curved stiffener, an arc butt plate \\
\bottomrule
\end{tabular}
\end{table}

\paragraph{Joint types.} Five plate joints are generated, named after ISO 17659: T (a
T-joint near $90^\circ$, an angle joint otherwise), corner, butt, lap and edge. A joint is
sampled from its type's ranges (Table~\ref{tab:ws_plate_ranges}): the thicknesses, the part
dimensions across the seam (they vary across it, never along it), the included angle, and
three defects with ISO 5817 limits, namely the root gap $g$, the linear misalignment $h$
and the angular misalignment $\beta$. A quality level (B, C, D or below D) is drawn first,
the defects are drawn within that level's limits, and the realised level is the worst of
the three.
\begin{table}[htbp]
\caption{Sampled plate-joint parameters (as coded).}
\label{tab:ws_plate_ranges}
\centering\small
\begin{tabular}{@{}L{0.27\textwidth}L{0.65\textwidth}@{}}
\toprule
Parameter & Range \\
\midrule
Thickness $t$ & U(1, 12) mm, the second part equal with $p=0.7$ (edge joints: U(1, 2) mm, ISO 9692-1) \\
Length along the seam & U(80, 400) mm per part, offset up to $\pm0.4\min L$ \\
Width across the seam & U(50, 250) mm \\
Included angle $\alpha$ & U(60, 120)$^\circ$ for T and corner; 180$^\circ$ butt; 0 lap and edge \\
Lap overlap & U(0.15, 0.60) of the upper plate's width \\
Root gap $g$ & U(0, min(ISO 5817 no.~617 limit, 2 mm)) for levels B--D; below D up to $\approx5.3$ mm \\
Misalignments & $h$ by ISO 5817 no.~5071; $\beta$ by level; $\beta\le0.4^\circ$ on lap and edge \\
Quality target & B, C, D, below D with 25\,\% each \\
\bottomrule
\end{tabular}
\end{table}
\TODO{Correct or state (weldset\_code\_review, P1--P2): parts interpenetrate on acute T
(28\,\% of specs, median 1 mm) and on edge/lap with $\beta>0$; $h$ and $\beta$ are stored but
not built on lap, edge and grooved butt; the quality level applies no.~617 to every type.}

\paragraph{Grooved butts.} Butt joints of thicker plates may carry an ISO 9692-1 groove,
chosen from the rows that the thickness admits: single-bevel (row 1.9.1), single-V (rows 1.3
and 1.5) and single-U (rows 1.6/2.6), with their bevel angle, root face, root gap and radius
drawn within the row.
\TODO{The row table (thickness windows, angles, root faces); the selection-thickness defect
(P1 \#5).}

\paragraph{Anti-shortcut construction.} A method can score well on synthetic plates for the
wrong reason: if every plate edge is parallel to a world axis, a detector can find the seam
from the axes alone. Two measures prevent this. Each assembly is turned by a random yaw over
the full circle, and plates get random convex outlines, so that only the seam-bearing edge
is constrained. A corpus gate then checks the result: the length-weighted histogram of angles
between free plate edges and the seam may not have more than twice the uniform mass in its
two terminal bins (threshold $2\cdot 20/90=0.444$).
\TODO{Gate history: 0.45 FAIL on the first Phase 4 corpus (rectangular grooved plates), 0.34
PASS after the fix (2026-09-07); the open swept-path seams report 0.45 (not gated).}

\paragraph{The fixture.} Half of the scenes in a fixture study stand on a fixture plate, a
slab of about $600\times400\times10$ mm tilted by up to $10^\circ$, so that a method meets
the plate-to-fixture contacts that a real cell has. Such contacts are never seams: a face
pair that involves the fixture is rejected. Fixture scenes are generated as twins of the
same seeds without the fixture.
\TODO{Plate twins are not at the same relative pose (P1 \#4).}
\SRC{weldgen/geom.py, joints.py, layouts.py, config.py, scene.py; PARAMETERS.md \S2--3.}

% -------------------------------------------------------------------------------------
\section{Curved Seams by Construction}
\label{sec:ws_curves}
For curved joints the order is reversed once more. A seam family is a curve family with a
condition under which two parts can realise it: the curve is drawn first and the parts are
derived from it, so the seam is exact by construction and no surface intersection is solved
at generation time. Seven families are generated: a line, a circle, an ellipse, a saddle, a
rounded rectangle, a swept path and an arc butt (Table~\ref{tab:ws_families}).
\begin{table}[htbp]
\caption{Curved seam families (as coded).}
\label{tab:ws_families}
\centering\small
\begin{tabular}{@{}L{0.17\textwidth}L{0.42\textwidth}L{0.33\textwidth}@{}}
\toprule
Family & Curve & Parts and draws \\
\midrule
circle & $\vect{c}_0+r(\cos\theta\,\hat{\vect{x}}+\sin\theta\,\hat{\vect{y}})$ & pipe of outer radius $r\sim$U(30, 80) on a plate \\
ellipse & $\vect{c}_0+a\cos t\,\hat{\vect{u}}+b\sin t\,\hat{\vect{v}}$, $a=b/\cos\tau$ & pipe of radius $b\sim$U(30, 70) tilted by $\tau$, mitred by the plate \\
saddle & intersection of two cylinders (Eq.~\ref{eq:ws_saddle}) & branch $r\sim$U(max(15, 0.3$R$), 0.7$R$) on a main pipe $R\sim$U(40, 80), tilt U(0, 30$^\circ$), offset $\pm0.4(R-r)$ \\
rounded rectangle & four straight runs and four quarter arcs of radius $r_c$ & box section $w,h\sim$U(60, 160) mm, $r_c\sim$U(8, $0.25\min(w,h)$) \\
swept path & a circular arc or a clamped cubic B-spline & curved stiffener of thickness U(4, 10), height U(50, 160) mm; the welds are the offsets at $\pm t/2$ \\
arc butt & an arc of radius $R\sim$U(100, 400) & two flat plates whose edges are the offset arcs at $R\mp g/2$ \\
\bottomrule
\end{tabular}
\end{table}

\figplaceholder{fig:ws_families}{One rendered scene per seam family (line, circle, ellipse, saddle,
rounded rectangle, swept path, arc butt) with the analytic seam drawn on it}{The seam families of
WeldSet.}

The saddle is the curve where a branch of radius $r$, axis $\hat{\vect{a}}$, meets a main pipe
of radius $R$, axis $\hat{\vect{m}}$ through $\vect{p}_m$. Around the branch circle
$\vect{c}(\varphi)=\vect{o}+r(\cos\varphi\,\hat{\vect{u}}+\sin\varphi\,\hat{\vect{v}})$, the point
$\vect{P}=\vect{c}+s\,\hat{\vect{a}}$ lies on the main cylinder when
\begin{equation}
s(\varphi)=\frac{-B+\sigma\sqrt{B^2-4AC}}{2A},\qquad
\begin{aligned}
A&=1-(\hat{\vect{a}}^{\top}\hat{\vect{m}})^2,\qquad \vect{d}=\vect{c}(\varphi)-\vect{p}_m,\\
B&=2\big[\vect{d}^{\top}\hat{\vect{a}}-(\vect{d}^{\top}\hat{\vect{m}})(\hat{\vect{a}}^{\top}\hat{\vect{m}})\big],\\
C&=\norm{\vect{d}}^2-(\vect{d}^{\top}\hat{\vect{m}})^2-R^2,
\end{aligned}
\label{eq:ws_saddle}
\end{equation}
with $\sigma=-1$ for the entering intersection; for an untilted, centred branch it reduces to
the textbook $z=\sqrt{R^2-r^2\sin^2\theta}$. The swept-path welds are offset curves of the
spine, $\vect{P}_o(t)=\vect{P}(t)+o\,\hat{\vect{n}}(t)$ with $\hat{\vect{n}}=\hat{\vect{z}}\times\hat{\vect{T}}$
and $o=\pm t/2$, because the spine itself lies in the middle of the material and is not a
weld.

\paragraph{Exactness checks.} Every curve has analytic positions; tangents are analytic except
on the B-spline, where they are central differences. Each part's ray-casting mesh records
its chord error, which must stay under 0.25 mm. At generation, a separate routine recomputes
the circle, ellipse and saddle from the placed parts by surface intersection, and the scene
is refused if the two differ by more than $10^{-6}$ mm; the measured differences are about
$10^{-13}$ mm.

\paragraph{Bores.} A pipe on a plate also has an inner fillet at its bore. A torch cone can
clear a 50 mm bore locally, but the torch body cannot enter it, so bores narrower than
80 mm are rejected as \emph{confined}; wider bores are weldable inner fillets.
\TODO{Corpus coverage differs from the draw (P2 \#9): the ellipse tilt is drawn
U(10$^\circ$, 40$^\circ$) but the corpus holds 10.2--30.6$^\circ$, because a dihedral under
about 59$^\circ$ fails the torch cone.}
\SRC{weldgen/curves.py, d29.py, constructors.py, verify\_curved.py, scene\_curved.py;
notes/patch\_phase6b.md (D33, D34).}

% -------------------------------------------------------------------------------------
\section{Seams by the Accessibility Rule}
\label{sec:ws_d4}
A seam is not every place where two parts touch: it is where they touch \emph{and} a torch
can reach. The accessibility rule (Algorithm~\ref{alg:ws_access}) turns this into a test on
every pair of faces of every pair of parts.

\begin{algorithm}[htbp]
\caption{The accessibility rule (plate joints)}
\label{alg:ws_access}
\begin{algorithmic}[1]
\REQUIRE placed parts with face registries; contact tolerance $\tau_c$
\ENSURE candidate seams, each weldable or rejected with a reason
\FORALL{face pairs $(F_A,F_B)$ of distinct parts}
  \IF{the normals are not parallel (more than $8^\circ$)}
    \STATE line $=\Pi_A\cap\Pi_B$ (Eq.~\ref{eq:plane_intersection}), clipped to both faces with slack $\tau_c$
    \STATE drop the pair if the faces are not mutually visible or are farther apart than $\tau_c$
  \ELSIF{the normals agree (coplanar faces)}
    \STATE line = the centreline of the in-plane gap, if the gap is within the tolerance
  \ELSE
    \STATE facing faces: reject (\emph{bisector blocked}); no line
  \ENDIF
  \STATE reject if a part is the fixture, if the dihedral is outside $[30^\circ,170^\circ]$, if the line is shorter than 10 mm
  \STATE search a torch axis (Eq.~\ref{eq:ws_axes}) whose cone clears the scene at five points; none: reject (\emph{bisector blocked})
\ENDFOR
\STATE demote a fillet that runs along a weldable butt centreline (\emph{toe of centreline}); demote cross-runs against the joint direction
\STATE class each seam from its faces; \emph{primary} = weldable $\wedge$ of the joint type's class
\end{algorithmic}
\end{algorithm}

\figplaceholder{fig:ws_access}{The accessibility rule on one T-joint: the three arms (intersecting,
facing, coplanar), the torch cone with its rim rays and probe depths, a cleared and a blocked
nook}{The accessibility rule.}

\paragraph{The torch cone.} The torch may tilt off the dihedral bisector $\hat{\vect{b}}$ in
steps of $15^\circ$ up to a work angle of $45^\circ$. The candidate axes are, nearest the
bisector first,
\begin{equation}
\hat{\vect{a}}=\cos\vartheta\,\hat{\vect{b}}+\sin\vartheta\,(\cos\phi\,\vect{u}+\sin\phi\,\vect{v}),\qquad
\vartheta\in\{0^\circ,15^\circ,30^\circ,45^\circ\},\ \phi=\tfrac{2\pi k}{8},
\label{eq:ws_axes}
\end{equation}
and an axis is admissible only if it has a positive component along both face normals. An
axis clears when a cone of $30^\circ$ half-angle (the axis and eight rim rays, probed at 3.75,
7.5 and 15 mm) is free of every part at five points along the seam; the first clearing axis
is stored as the seam's approach direction. A consequence worth stating: an open nook clears
only if its dihedral exceeds about $60^\circ$, so the effective floor is $60^\circ$, not the
nominal $30^\circ$.

\paragraph{Classes and the three seam curves.} A seam's class follows mechanically from its
faces: two broad faces give a fillet, or a butt above $172^\circ$; two edge faces give an
edge seam; one of each gives a lap toe. With a root gap there are three natural curves. The
\emph{nominal} curve is the intersection of the extended face planes; the \emph{root} lies
on the terminating face; the \emph{gap-mid} lies between the faces. The nominal curve is
stored and scored, because it is continuous as $g\to0$ and does not depend on which part is
called A; the other two are stored alongside for plate seams. For a fillet whose terminating
edge stands $H_e$ above the base plane, the root lies $\delta_{\mathrm{root}}=H_e/\sin\varphi$
from the nominal line.
\TODO{Fix the coplanar-arm bug (P1 \#3); the three gap-mid definitions (P2 \#10); which
rejected candidates are stored (P2 \#13).}
\SRC{weldgen/accessibility.py; SCHEMA.md \S1.3, \S2.6; MANUAL.md \S3.3.}

% -------------------------------------------------------------------------------------
\section{Tack Placement Rule}
\label{sec:ws_tacks}
Tacks are a rule applied to the seams, versioned and re-runnable, not a label: a tack layout
is a shop convention, and someone who disagrees with its constants can re-run the rule. Only
primary seams are tacked ("you tack what you weld"). With $t$ the gauge of the thinnest
workpiece in the scene (for a swept band, the smaller of its thickness and height),
tackrule-0.1 sets the tack
length, the spacing bounds and the end margin
\begin{equation}
\begin{gathered}
\ell = \operatorname{clip}(4t,\,10,\,50)\ \mathrm{mm},\qquad m = \max(2t,\ \ell),\\
d_{\min} = 10t,\qquad d_{\max} = \min(33t,\,400)\ \mathrm{mm},
\end{gathered}
\label{eq:tackrule_len}
\end{equation}
and places the tacks by Algorithm~\ref{alg:ws_tacks}. For an open seam of length $L$ with
$L_e=L-2m$, the count before the lower bound is
\begin{equation}
n = \max\!\left(2,\ \ceil*{\frac{L_e}{d_{\max}}}+1\right),
\label{eq:tackrule_n}
\end{equation}
and the tacks sit at $s_k=m+k\,L_e/(n-1)$, so that both ends are always tacked.

\begin{algorithm}[htbp]
\caption{tackrule-0.1 on one primary seam}
\label{alg:ws_tacks}
\begin{algorithmic}[1]
\REQUIRE seam of length $L$, open or closed; thinnest member gauge $t$
\STATE compute $\ell, d_{\max}, d_{\min}, m$ (Eq.~\ref{eq:tackrule_len})
\IF{the seam is open}
  \IF{$L-2m\le\ell$}
    \STATE one tack at $L/2$ (recorded as degenerate) and \textbf{return}
  \ENDIF
  \STATE $n$ from Eq.~\ref{eq:tackrule_n}; \textbf{if} $n>2$ and $L_e/(n-1)<d_{\min}$ \textbf{then} $n\gets\max(2,\lfloor L_e/d_{\min}\rfloor+1)$
  \STATE $s_k=m+kL_e/(n-1)$, $k=0,\dots,n-1$; order: the ends, then by recursive bisection
\ELSE
  \STATE $n=\max(4,\lceil L/d_{\max}\rceil)$ rounded up to even; \textbf{if} $L/n<d_{\min}$ \textbf{then} $n\gets\max(4,2\lfloor L/(2d_{\min})\rfloor)$
  \STATE phase $\phi$ from the first eight bytes of sha256(scene, seam) / $2^{64}$
  \STATE $s_k=(\phi L+kL/n)\bmod L$; order: opposite pairs in bisection order
\ENDIF
\STATE interleave seams of the same class round-robin in the weld order
\end{algorithmic}
\end{algorithm}

\figplaceholder{fig:ws_tacks}{Tack layouts of tackrule-0.1 on an open seam (both ends, bisection
order) and on a closed loop (even count, hashed phase, opposite pairs)}{Tack layouts on an open and a
closed seam.}

The phase of a closed loop is taken from a hash so that it is deterministic but not tied to
the curve's parametrisation; any score on a closed seam must therefore be invariant to
rotation along the loop. The constants are shop conventions: the thickness-scaled spacing
and the 400 mm ceiling follow JASS 6 (through the Kobelco handbook), the tack length is
attributed to EN 1011-2, and Tomk\'ow et al.\ \cite{tomkow2020tacks} give the literature
anchor.
\TODO{Obtain the JASS 6 / EN 1011-2 texts before citing (the length rule is attributed, not
verified against the standard).}
\SRC{weldgen/tacks.py; MANUAL.md \S10; scripts/apply\_rule\_blocks.py.}

% -------------------------------------------------------------------------------------
\section{The Most Probable Seam}
\label{sec:ws_mps}
A second task asks a question that a single view can answer: given what one camera sees, at
which seam is a coarsely positioned torch? The label is again a versioned rule
(mps\_rule-0.1), a geometric proxy by decision: the weldable seam with the largest visible
arc length, ties broken by the larger fold $|180^\circ-\text{dihedral}|$, then by the lower
identifier, and no seam if the largest visible length is under 10 mm. The margin between the
runner-up and the winner, $\ell_{(2)}/\ell_{(1)}\in[0,1]$, says how much of a choice the view
leaves.
\SRC{weldgen/mps.py; MANUAL.md \S1.4, \S10.}

% -------------------------------------------------------------------------------------
\section{Sensor Model, Visibility and Rendering}
\label{sec:ws_sensor}

\paragraph{Camera.} Each scene has a pinhole camera of $1280\times720$ pixels with the
principal point at the image centre and one of three stereo profiles
(Table~\ref{tab:ws_profiles}). The camera looks at an aim point near the joint,
$\vect{a}=\vect{m}+0.15\,L\,\boldsymbol{\xi}$ with $\vect{m}$ the mean of the weldable seam
midpoints, $L$ the longest part edge and $\boldsymbol{\xi}\sim\mathrm{U}([-1,1]^3)$, from
\begin{equation}
\begin{gathered}
\vect{C}=\vect{a}+s\,(\cos e\cos\alpha,\ \cos e\sin\alpha,\ \sin e)^{\top},\\
s=\operatorname{clip}\!\Big(\frac{f\,E}{\varphi\,\min(W,H)},\ \max(z_{\min},300),\ 1200\Big)\ \mathrm{mm},
\end{gathered}
\label{eq:ws_pose}
\end{equation}
with elevation $e\sim$U(15$^\circ$, 85$^\circ$), azimuth $\alpha\sim$U(0, 360$^\circ$), a roll
of up to $\pm15^\circ$, $E$ the longest side of the assembly's bounding box and a framing
fraction $\varphi\sim$U(0.35, 1.45).
\begin{table}[htbp]
\caption{Sensor profiles (as coded).}
\label{tab:ws_profiles}
\centering\small
\begin{tabular}{@{}lcccc@{}}
\toprule
Profile & baseline $b$ (mm) & focal length $f$ (px) & disparity noise $\sigma_d$ (px) & min.\ depth (mm) \\
\midrule
d435i & 50 & 674 & 0.08 & 280 \\
stereo\_good & 120 & 1100 & 0.05 & 200 \\
stereo\_poor & 35 & 450 & 0.15 & 400 \\
\bottomrule
\end{tabular}
\end{table}

\paragraph{Sampling and visibility.} Every face is sampled area-uniformly at a density drawn
from U(0.25, 4) points per mm$^2$, with exact positions and normals. A point is visible from
the camera when it is inside the image and beyond the minimum depth, faces the camera,
$(\vect{C}-\vect{p})^{\top}\vect{n}>0$, and the ray to the camera hits no part. Rays are
tested exactly against planar parts (slab and convex-prism clipping) and by an entry-only
Möller--Trumbore test against the 0.05 mm chord meshes of tubes and swept parts. A second
flag marks a point \emph{exterior} when a scan from outside could reach it: some ray of a fan
within $70^\circ$ of its normal (the normal itself and eight rays each at $38.5^\circ$ and
$66.5^\circ$) escapes every part. The exterior points form the input of
the full-scan condition, and the visible points that of the single-view condition.
\TODO{Fix and re-run before final numbers (P1 \#1): the exterior flag is computed on a 2 mm
voxel subsample and is wrong on thin stacked plates (edge 12\,\%, lap 2.4\,\% of points).
Grooved plates use an approximate (9-sample) occlusion test.}

\paragraph{Noise.} The cloud is stored noiseless, and noise is applied at evaluation time by
a released function, so that the realisation is a rule, not data. A stereo camera measures
depth from disparity, $z=fb/d_{\mathrm{px}}$, so a disparity error $\sigma_d$ gives
\begin{equation}
\sigma_z(z)=\frac{\sigma_d\,z^{2}}{f\,b},\qquad
\tilde{\vect{X}}=\vect{X}+\epsilon_0\,\sigma_z(z)\,\hat{\vect{r}}+\sigma_\ell(z)\,(\epsilon_1\hat{\vect{x}}_c+\epsilon_2\hat{\vect{y}}_c),\qquad
\sigma_\ell=0.8\,z/f,
\label{eq:ws_noise}
\end{equation}
with the depth term along the viewing ray $\hat{\vect{r}}$, a lateral term of 0.8 pixel along
the image axes, $\epsilon_i\sim\mathcal{N}(0,1)$, and points dropped beyond $75^\circ$ of
incidence or inside the minimum depth. For the d435i profile this gives 0.59 mm at 0.5 m and
2.37 mm at 1 m. The benchmark scales $(\sigma_d,\sigma_\ell)$ by $k\in\{0,1,2\}$.

\figplaceholder{fig:ws_sensor}{(a) $\sigma_z$ against range for the three profiles
(Equation~\ref{eq:ws_noise}); (b) one scene: the exact cloud, the visible points of the single view,
the exterior points, and the noisy cloud at $k=1$}{The sensor model.}

\paragraph{Rendered tier.} The second tier renders scenes of the same generator in Isaac Sim
(RGB, depth and seam, tack and object masks) with published material reflectances, real
backgrounds and lighting. A twin gate ties the two tiers together: the rendered clean depth
of the first view, back-projected, must land on the exact primitives (99th-percentile residual
under 0.25 mm) and agree with the analytic visibility (coverage $\ge0.90$, object identity
$\ge0.98$). The gate also served as an independent check of the first tier: it found three
ray-test defects in the analytic visibility (sampled containment stepping through thin tube
walls, a mesh cache keyed before re-posing, and exit crossings counted as occluders), which
were fixed and the benchmark regenerated before its results were computed.
\TODO{State precisely (C\_weldset \S2.7): the rendered corpus is train\_v1 (3600 scenes $\times$
10 views, its own seeds), not the benchmark corpus; one material reflectance (bronze) is
interpolated; mask widths are nominal. Gate statistics per stratum: DESKTOP.}
\SRC{weldgen/camera.py, sampling.py, visibility.py, noise.py, weldgen/render/;
PARAMETERS.md \S4; phase8\_plan.md.}

% -------------------------------------------------------------------------------------
\section{The Corpus and Its Strata}
\label{sec:ws_corpus}
The benchmark corpus has twelve strata of 60 scenes each: six T strata (a line and the five
curved T families), three butt strata (square, grooved and arc), and the corner, lap and
edge lines. Each stratum is built from one configuration over a fixed seed range, and the
strata are reported separately. Table~\ref{tab:ws_strata} gives the ranges.
\begin{table}[htbp]
\caption{WeldSet strata and their geometry (as coded; emitted where it differs).}
\label{tab:ws_strata}
\centering\small
\begin{tabular}{@{}L{0.2\textwidth}L{0.72\textwidth}@{}}
\toprule
Stratum & Geometry \\
\midrule
T/line & plates $t$ 1--12 mm, $\alpha$ 60--120$^\circ$, gap 0--5.3 mm; yaw over the full circle, outline on the stem \\
corner/line & as T, outlines on both parts \\
butt/line (square) & $t$ 1--12 mm, $\alpha=180^\circ$, flush on one side when thicknesses differ \\
butt/line (grooved) & $t$ 4--20 mm; single-bevel, single-V or single-U by ISO 9692-1 row \\
butt/arc & arc radius 100--400 mm, gap 0.3--2.0 mm, plates $t$ 3--10 mm \\
lap/line & $t$ 1--12 mm, overlap 15--60\,\% of the upper plate, $\beta\le0.4^\circ$ \\
edge/line & $t$ 1--2 mm \\
T/circle & pipe radius 30--80 mm, wall 3--8 mm, base plate 4--12 mm \\
T/ellipse & pipe radius 30--70 mm, tilt drawn 10--40$^\circ$ (emitted 10--31$^\circ$) \\
T/saddle & main pipe 40--80 mm, branch 0.3--0.7 of it (at least 15 mm), tilt 0--30$^\circ$ \\
T/rounded rect & box section 60--160 mm, corner radius 8 mm up to a quarter of the short side \\
T/swept path & stiffener 4--10 mm thick, 50--160 mm high, span 120--300 mm \\
\bottomrule
\end{tabular}
\end{table}
\TODO{Corpus sizes and yields per stratum, the fixture twins (634: 274 plate + 360 curved) and
train\_v1: DESKTOP; the final corpus is the visibility-fixed rebuild of 2026-09-12.}
\SRC{scripts/make\_bench6b.py; configs/; review C2 (ranges as coded).}

% -------------------------------------------------------------------------------------
\section{Evaluation Protocol}
\label{sec:ws_protocol}

\paragraph{Tasks.} Task~1 asks for every seam from a perfect multi-view exterior scan of the
assembly; Task~2 asks for the most probable seam (\S\ref{sec:ws_mps}) from a single view.

\paragraph{Metrics.} A predicted polyline set $\mathcal{P}$ and the truth $\mathcal{G}$ (the
nominal curves of the primary seams) are densified at 1 mm, and a predicted point counts if
its exact distance to a truth segment is within $\tau=3$ mm:
\begin{equation}
\begin{gathered}
\mathrm{prec}=\frac{1}{|\tilde{\mathcal{P}}|}\sum_{\vect{x}\in\tilde{\mathcal{P}}}\mathbf{1}\big[d(\vect{x},\mathcal{G})\le\tau\big],\qquad
\mathrm{rec}=\frac{1}{|\tilde{\mathcal{G}}|}\sum_{\vect{y}\in\tilde{\mathcal{G}}}\mathbf{1}\Big[\min_{\vect{x}\in\tilde{\mathcal{P}}}\norm{\vect{y}-\vect{x}}\le\tau\Big],\\
F_1=\frac{2\,\mathrm{prec}\cdot\mathrm{rec}}{\mathrm{prec}+\mathrm{rec}},
\end{gathered}
\label{eq:ws_f1}
\end{equation}
where $\tilde{\mathcal{P}}$ and $\tilde{\mathcal{G}}$ are the densified point sets and
$d(\vect{x},\mathcal{G})$ is the exact distance to the nearest truth segment.
The literature's own path metric is reported as well: predictions are matched one-to-one to
truth seams, greedily by mean distance, and each matched path, resampled at 0.25 mm, gives
the root-mean-square and the maximum of its distances to the truth.
\TODO{Chamfer and the lateral distribution are computed but not reported; matching has no
distance cap; Task~1 scores primary seams only (weldset\_code\_review, ICRA list).}

\paragraph{The oracle ladder.} Most published pipelines start with a learned coarse stage
(a segmentation network, a detector, a keypoint model) that was not released. Re-implementing
the geometric stage alone would handicap each method by a stage its authors never published.
At level L0, therefore, each method receives \emph{its own paper's} coarse stage from the
ground truth (Table~\ref{tab:ws_methods}); level L1 withholds it and changes nothing else.
The difference between the two is the price of the unpublished stage.

\begin{table}[htbp]
\caption{The seven re-implemented methods and their L0 oracles.}
\label{tab:ws_methods}
\centering\small
\begin{tabular}{@{}L{0.07\textwidth}L{0.42\textwidth}L{0.43\textwidth}@{}}
\toprule
 & Method & L0: the coarse stage given from truth \\
\midrule
M1 & multi-plane RANSAC, pairwise plane intersection \cite{yi2026seam} & a 40 mm band around the seams (their PointNet++) \\
M2 & curvature-seeded region growing, seam points between two grown surfaces \cite{wei_regiongrow} & a keypoint crop keeping 12.7\,\% of the cloud plus per-surface labels (their FastSAM) \\
M3 & local oriented bounding-box features, crease points, curve fit \cite{zhang_lobb,zhang2025lobbdesc} & per-part labels (their K-Net) \\
M4 & plane extraction by point-pair coplanarity, orthogonal-pair voting \cite{wang2024ppfseam} & the 40 mm seam band (their Faster R-CNN box) \\
M5 & PCA slicing of the weld cloud, centre per slice, B-spline \cite{wang2026pipes} & one mask per seam (their YOLO + DeepLab) \\
M6 & CAD model registered by coherent point drift, seam transferred \cite{fang2024modelreg} & the model itself (no L1) \\
M7 & plane or quadric per welding surface, seam near the intersection \cite{li2026quadric} & per-surface labels and part membership (their region growing) \\
\bottomrule
\end{tabular}
\end{table}

\paragraph{Fidelity.} Each method has its own notebook recording every reading of its paper,
every ambiguity and, where the paper reports one, a reproduction of its number. On
2026-09-14 every module was re-read against its paper; six readings were corrected and the
old readings kept as labelled rungs of the ladder, so that each choice is priced. The first
author of M3 reviewed the implementation choices and endorsed them in writing.
\TODO{One paragraph per method in Chapter 2 (mechanism, coarse stage, the reading choices);
the reproductions (M1: RMSE 0.60 mm / max 0.81 mm against 0.64 / 1.16 reported).}

\paragraph{Conditions.} The two views (exterior scan, single view); L0 and L1 (L1 on the
exterior scan); noise at $k=1,2$ on the single view; the fixture twins (40 plate pairs per
method); 30 seeds for the randomised M1. The harness was validated on a fake oracle (the
truth plus known jitter) before any method ran: a jitter of $\sigma$ returns an error of about
$\sigma$, a phantom seam costs precision and never steals a match, and a miss is an unmatched
row.
\SRC{scripts/baselines/metrics.py, harness.py, lit\_*.py; run\_phase4\_batch.py; notebooks
04--11, 15; review C4, C5, C6.}

% -------------------------------------------------------------------------------------
\section{Proposed Seam Extractors}
\label{sec:ws_own}
Two extractors were developed in this work and scored on the same corpus with the same
harness and metrics as the seven literature methods. Neither needs a CAD model, and each
starts from an input that the robotic cell produces anyway (\S\ref{sec:ws_link}).
\TODO{Attribution: the footprint idea is İsmail Hoca's; say so (here or in the
acknowledgments).}

\subsection{Footprint Extraction from Sequential Scans}
\label{sec:ws_footprint}
The idea is to scan the base part $A$ alone, then the assembly $A+B$. Whatever the first scan
does not explain belongs to $B$, so the parts are separated by the order of placement instead
of by a learned model. Where $B$ rests on $A$, each part leaves a \emph{footprint} on the
other, and the seam lies at its edge. All thresholds are multiples of the voxel size $v$, so
the method scales with the data. The distance from a point $\vect{x}$ to a sampled surface
$Q$ with normals $\vect{n}_i$ is bounded so that a coplanar neighbour (a butt joint) does not
read zero everywhere:
\begin{equation}
D_Q(\vect{x})=\min_{i\in\mathcal{N}_4(\vect{x})}\sqrt{\big((\vect{x}-\vect{q}_i)^{\top}\vect{n}_i\big)^2+\max\!\big(\norm{(\vect{x}-\vect{q}_i)_\perp}-r_s,\,0\big)^2},\qquad r_s=0.75v,
\label{eq:fp_distance}
\end{equation}
and the footprints are the points of each part closer than $g_{fp}=3v+g$ to the other,
\begin{equation}
F_A=\{\vect{p}\in A: D_B(\vect{p})<g_{fp}\},\qquad F_B=\{\vect{q}\in B: D_A(\vect{q})<g_{fp}\},
\label{eq:footprint}
\end{equation}
where $g$ is the expected fit-up gap of the welding procedure. Algorithm~\ref{alg:ws_footprint}
gives the steps.

\begin{algorithm}[htbp]
\caption{Footprint seam extraction}
\label{alg:ws_footprint}
\begin{algorithmic}[1]
\REQUIRE base scan $A_0$, assembly scan $C$, expected gap $g$
\STATE voxel-downsample both at $v$ (the median 4th-neighbour distance); estimate normals
\STATE optionally register $A_0$ into $C$ (point-to-plane ICP)
\STATE split $C$: a point near $A_0$ (and with a compatible normal) belongs to $A$, the rest to $B$
\STATE compute $D_B$ on $A$ and $D_A$ on $B$ (Eq.~\ref{eq:fp_distance}); the footprints $F_A$, $F_B$ (Eq.~\ref{eq:footprint})
\FORALL{footprint points}
  \STATE slide along the negative surface gradient of $D$, re-projected onto the point's own surface, until $D=g$ (secant iterations)
\ENDFOR
\STATE keep the slid points with both parts in view within $3v$ (exposure)
\STATE cluster them by position and opening direction; thin, order by the longest path of a spanning tree, fit a smoothing spline, resample
\STATE per path point: leg directions $\to$ joint angle $\theta$, bisector, torch axis; roll a ball of radius $r$ in the cross-section $\to$ toes, legs, bead area
\STATE split paths at sharp turns; classify points as butt, open corner, fillet or narrow by $\theta$
\end{algorithmic}
\end{algorithm}

\paragraph{The rolling ball.} Along the bisector $\hat{\vect{b}}$ from a root point
$\vect{x}$, the ball centre $\vect{c}=\vect{x}+h\hat{\vect{b}}$ is the first $h$ at which
$\min(d_A,d_B)\ge r$; for an ideal fillet of opening $\theta$ this gives
$h=r/\sin(\theta/2)$ and legs $\ell=r\cot(\theta/2)$.
\TODO{How WeldSet runs it: the base scan is synthesised (an independent sample of part A with
A as the only occluder); registration is off (exact poses); the gap given is the constructed
root gap ("WPS" arm), with a zero-gap rung. Limitations: grooved butts (no footprint on the
top faces), thin laps and edges, closed rings recovered as partial arcs, the rolling-ball
outputs not scored.}
\SRC{Self\_ideas/test1\_ismail\_hoca/weld\_seam.py; scripts/baselines/self\_footprint.py;
notebooks/build\_16; review C\_weldset \S3.A.}

\subsection{Matching and Sphere-Tracing on Labelled Clouds}
\label{sec:ws_seamfind}
Given two point clouds that are already labelled by part, a point of $A$ whose nearest
point on $B$ lies within the gap bound is near the seam. From there, each side can walk
towards the other over its own surface until the two meet: the place where the walkers stop
gives the toes, and the toes give the gap and the root. Normals are estimated per part, so
they do not blur across the joint, and each part gets a moving-least-squares projection
$\Pi_A$, $\Pi_B$ onto its local surface patch. With $\vect{q}_k$ the closest point of $B$'s
finite patch and $\vect{n}_k$ $A$'s normal, a walker on $A$ takes the steps
\begin{equation}
\vect{x}_{k+1}=\Pi_A\big(\vect{x}_k+\alpha\,(\mathbf{I}-\vect{n}_k\vect{n}_k^{\top})(\vect{q}_k-\vect{x}_k)\big),\qquad
d_k=\norm{\vect{x}_k-\vect{q}_k},\qquad \alpha=0.9,
\label{eq:sphere_trace}
\end{equation}
i.e.\ $\alpha$ times the tangential part of the vector to $B$, so that it cannot overshoot.
It stops when it touches $B$ ($d_k<0.05h$), stalls at a gap ($d_{k-1}-d_k<0.01h$), or reaches
the edge of its own patch. On two planar faces meeting at an angle $\theta$, each step shrinks
the distance by the same ratio, $d_{k+1}=d_k(1-\alpha\sin\theta)$, so the mean step ratio
$\bar r$ gives a second, independent estimate of the joint angle,
$\theta_r=\arcsin\big((1-\bar r)/\alpha\big)$.

\begin{algorithm}[htbp]
\caption{Matching and sphere-tracing seam extraction}
\label{alg:ws_seamfind}
\begin{algorithmic}[1]
\REQUIRE labelled clouds $A$, $B$ (and the camera position for normal signs); gap bound $g_{\max}$
\STATE per part: downsample at its spacing $h$; normals (flattest neighbourhood); moving-least-squares projection $\Pi$
\STATE candidates: points within $\tau=g_{\max}+2h$ of the other part; drop those no torch could see (cone prefilter)
\STATE seeds: pair each candidate with its nearest candidate on the other part; add coplanar boundary pairs (butts)
\STATE class per seed by the normal angle: fillet, coplanar, facing (dropped), parallel (dropped)
\FORALL{seeds}
  \STATE walk on $A$ towards $B$ and on $B$ towards $A$ (Eq.~\ref{eq:sphere_trace}) $\to$ toes $\vect{T}_A$, $\vect{T}_B$; gap $=\norm{\vect{T}_A-\vect{T}_B}$
  \STATE refit each joint face's plane away from the crease; root = the point of the two planes' intersection line nearest the seed (fillet) or the toes' midpoint (coplanar)
\ENDFOR
\STATE cluster the roots (by position, then by the normal pair); order by the longest path of a spanning tree; smooth and resample at 2 mm
\STATE apply the generator's rules: drop cross-runs and fillets along a butt centreline
\STATE torch check at every point: a 30$^\circ$ cone over 7 work $\times$ 3 travel angles, and the 80 mm bore rule; output the weldable spans
\end{algorithmic}
\end{algorithm}
\TODO{How WeldSet runs it: the labels are the exact part identifiers (the same input M3 gets at
L0); the gap bound is 5 mm; the defaults differ from the plan in 17 documented ways and were
chosen on a 12-scene smoke round of the same corpus (say so). Limitations: thin sheets, grooved
butts in single view, the gap estimate biased upward, $\theta_r$ undefined near 90$^\circ$.}
\SRC{Self\_ideas/test2\_cagdas/seam\_finder.md, seamfind/; scripts/baselines/self\_seamfind.py;
notebooks/build\_17; review C\_weldset \S3.B.}
\figplaceholder{fig:ws_own}{The two extractors on one T-joint: (a) base scan, assembly scan,
footprints and the slid points; (b) seeds, the two walkers' paths, toes and the root}{The two
proposed seam extractors.}

% -------------------------------------------------------------------------------------
\section{WeldSet and the Robotic Cell: One Geometry}
\label{sec:ws_link}
WeldSet and the cell share one geometry, in two directions.

\textbf{From WeldSet to the cell.} The cell represents each registered CAD part as a WeldSet
primitive (\S\ref{sec:mc_registry}) and applies the tack rule of \S\ref{sec:ws_tacks}
unchanged. Its seams follow WeldSet's geometry in a form adapted to registered rather than
constructed parts: the plate rule is the accessibility rule with a pose tolerance in place of
the contact tolerance, and curved seams are computed by intersecting the registered surfaces
with WeldSet's curve library (plane--cylinder, cylinder--cylinder, offset curves) and judged
by its curved-seam verdict. The seams and tacks the robot marks are therefore WeldSet's exact
geometry evaluated at the measured poses; their error is the registration's, which the pen's
contact measures independently (Eq.~\ref{eq:touch}).

\textbf{From the cell to WeldSet.} The two proposed extractors take as input what the cell's
procedure already produces (Table~\ref{tab:ws_link}). Both are CAD-free, so they could check
the cell's CAD-based seams, or replace them where no CAD exists; WeldSet scores them against
exact truth before they meet a real scan.
\begin{table}[htbp]
\caption{The extractors' inputs in WeldSet and in the cell.}
\label{tab:ws_link}
\centering\small
\begin{tabular}{@{}L{0.2\textwidth}L{0.33\textwidth}L{0.39\textwidth}@{}}
\toprule
Input & In WeldSet & In the cell \\
\midrule
Scan of the base alone & synthesised: an independent sample of part A & the parts are registered and saved one at a time, so the base is scanned before the next part is placed \\
Fit-up gap & the constructed root gap & a procedure's nominal gap, or the measured fit-up \\
Part labels & exact part identifiers & SAM2 masks and the registration label every point \\
\bottomrule
\end{tabular}
\end{table}
\TODO{DECIDE (THESIS\_PLAN \#8): run the two extractors on the cell's real scans (then a
results section) or keep WeldSet scores only and name the real-scan run as future work.}
\figplaceholder{fig:ws_link}{WeldSet and the cell: shared primitives and tack rule; the
seam rule in its registered form; sequential scans and part labels feeding the two
extractors}{How WeldSet and the robotic cell are connected.}

% -------------------------------------------------------------------------------------
\section{The Annotation-Error Experiment}
\label{sec:ws_annotation}
The analytic truth also allows the field's usual reference to be measured. Twenty scenes, four
per plate joint type, were exported as anonymous point clouds with no truth and labelled in
CloudCompare by a briefed engineering student following the literature's protocol, blind to
the truth (4--10 clicks per seam, 30 primary truth seams). The clicks were scored against the
analytic seams with a 15 mm acceptance gate, after the toolchain had been validated on a
synthetic annotator with known noise, one planted miss and one planted phantom. The result,
the label-noise floor of hand annotation, is reported in Chapter~\ref{chp:results}.
\SRC{scripts/annotation/; out/annotation (laptop copy; confirm on the desktop); ICRA \S V.}

% -------------------------------------------------------------------------------------
\section{Summary}
\label{sec:ws_summary}
\TODO{Synthesis paragraph (Kurt's chapter closer): what WeldSet provides (exact truth,
dials, a fair protocol) and how it connects to the cell.}
